% ===========================================================================
% Twin Decision Theory: group-theoretic discrimination for twin spaces
% Companion note to arithmetics/symbolic/twin/decision.py
% ===========================================================================
\documentclass[11pt,a4paper]{article}
\usepackage[utf8]{inputenc}
\usepackage[T1]{fontenc}
\usepackage{lmodern}
\usepackage[english]{babel}
\usepackage{amsmath,amssymb,amsthm,amsfonts}
\usepackage{booktabs}
\usepackage[margin=1in]{geometry}
\usepackage{xcolor}
\usepackage{hyperref}
\hypersetup{colorlinks=true, linkcolor=blue!60!black, urlcolor=blue!60!black}

\newcommand{\Op}{\mathrm{Op}}
\newcommand{\Z}{\mathbb{Z}}
\newcommand{\R}{\mathbb{R}}
\newcommand{\Inn}{\mathrm{Inn}}
\newcommand{\Aut}{\mathrm{Aut}}
\newcommand{\Fix}{\mathrm{Fix}}
\newcommand{\Stab}{\mathrm{Stab}}
\newcommand{\code}[1]{\texttt{#1}}
\DeclareMathOperator{\Norm}{N}

\theoremstyle{plain}
\newtheorem{theorem}{Theorem}
\newtheorem{proposition}[theorem]{Proposition}
\newtheorem{corollary}[theorem]{Corollary}
\newtheorem{lemma}[theorem]{Lemma}
\theoremstyle{definition}
\newtheorem{definition}[theorem]{Definition}
\theoremstyle{remark}
\newtheorem{remark}[theorem]{Remark}

\title{\textbf{Twin Decision Theory}\\[0.3em]
\large Group-theoretic discrimination for twin operation spaces}
\author{Jocelyn Nemb\'e \\
\small Modeling and Calculus Laboratory, Libreville \\
\small \texttt{jnembe@roots-insights.org}}
\date{June 2026}

\begin{document}
\maketitle

\begin{abstract}
\noindent
We give a rigorous, unified framework casting the selection of a twin operation
as a \emph{decision problem}, and we promote the structural invariants of group
theory --- conjugacy classes, centralizers, the centre, and the
\emph{class equation} --- into instruments of discrimination. We first correct
the classification statement of the foundations volume: the operation space
$\Op$ is canonically the homogeneous space $G/L$ (a coset space), and it carries
a group structure if and only if the structural kernel $L$ is \emph{normal}. We
exhibit a concrete counterexample ($S_4$ acting on $+\bmod 4$, where $L$ has
order $2$ and is not normal). We then prove the reduction principles: a cost
that is a class function is constant on conjugacy classes, so the class equation
\emph{enumerates the decision cells}; in general the cells are the orbits of the
normalizer $\Norm_G(L)$, counted by Burnside. Finally we record the compact-Lie
instance, where the class equation becomes the Weyl integration formula
(worked out for $SU(2)$). Every statement is matched to a computational test in
\code{arithmetics/symbolic/twin/decision.py}.
\end{abstract}

\section{The twin decision problem}

\begin{definition}[Twin decision problem]
A \emph{twin decision problem} is a tuple $\mathcal{D} = (E, \odot_e, G, \rho, J)$
where $(E,\odot_e)$ is a magma, $\rho : G \to \mathrm{Sym}(E)$ is a group action
(written $g\cdot x$), and for each $g\in G$ the \emph{twin operation} is
\[
x \odot_{eg} y := g^{-1}\cdot\bigl((g\cdot x)\odot_e(g\cdot y)\bigr).
\]
Let $\pi : G \to \Op(E,\odot_e;G)$, $g\mapsto \odot_{eg}$, be the
\emph{transport map} onto the set $\Op$ of distinct twin operations, and
\[
L := \{\,g\in G : \odot_{eg} = \odot_e\,\}
\]
the \emph{structural kernel}. Finally $J : \Op \to \R_{\ge 0}$ is a \emph{cost}.
To \emph{solve} $\mathcal{D}$ is to compute $\odot^\star = \arg\min_{\Op} J$.
\end{definition}

All earlier applications --- batched arithmetic (speed), dynamic range,
function approximation, optimal-base selection, twin differential calculus ---
are instances of $\mathcal{D}$ with different $G$ and $J$. This is the single
formulation sought.

\section{The kernel and the corrected classification}

\begin{proposition}[Kernel]\label{prop:kernel}
$L$ is a subgroup of $G$, equal to the automorphism group of $(E,\odot_e)$
inside $\rho(G)$.
\end{proposition}

\begin{proposition}[Fibres are right cosets]\label{prop:fibres}
For $g,h\in G$,
\[
\odot_{eg} = \odot_{eh} \iff h g^{-1}\in L \iff Lg = Lh.
\]
\end{proposition}

\begin{proof}
$\odot_{eg}=\odot_{eh}$ means
$g^{-1}\!\cdot\!\bigl((g\cdot x)\odot_e(g\cdot y)\bigr)
 = h^{-1}\!\cdot\!\bigl((h\cdot x)\odot_e(h\cdot y)\bigr)$ for all $x,y$. Apply
$h\cdot(-)$ and substitute $u=g\cdot x,\ v=g\cdot y$ (a bijection of $E$); with
$m := hg^{-1}$ this reads $m\cdot(u\odot_e v) = (m\cdot u)\odot_e(m\cdot v)$ for
all $u,v$, i.e. $m\in L$. Hence $hg^{-1}\in L$, equivalently $Lg=Lh$.
\end{proof}

\begin{corollary}[Classification]\label{cor:classification}
The transport map induces a bijection $L\backslash G \xrightarrow{\ \sim\ }\Op$,
$Lg \mapsto \odot_{eg}$. In particular $|\Op| = [G:L]$.
\end{corollary}

\begin{theorem}[Group structure $\iff$ normality]\label{thm:group-iff-normal}
The rule $\odot_{eg}\ast\odot_{eh} := \odot_{e,gh}$ is a well-defined binary
operation on $\Op$ (making $\pi$ a homomorphism and $\Op\cong G/L$ a group) if
and only if $L \trianglelefteq G$.
\end{theorem}

\begin{proof}
Replace $g,h$ by $ag, bh$ with $a,b\in L$. Then
$L\,(ag)(bh) = L\,g b h$, which equals $L\,gh$ for all $b\in L$ iff
$g b g^{-1}\in L$ for all $b\in L$ and all $g$, i.e. iff $g^{-1}Lg = L$ for all
$g$ --- precisely $L\trianglelefteq G$.
\end{proof}

\begin{remark}[Correction to the foundations volume]\label{rem:correction}
The foundations volume asserts that $\pi$ is always a group homomorphism and
hence $L\trianglelefteq G$. This is \emph{false in general}. For
$G = S_4$ acting on $E=\{0,1,2,3\}$ with base operation $+\bmod 4$, the kernel
$L$ has order $2$ and is \emph{not} normal (a transposition subgroup; $S_4$ has
no normal subgroup of order $2$). Thus $\Op$ is the homogeneous space $G/L$ with
$|\Op| = 12$, but is \emph{not} a group. The correct universal statement is
Corollary~\ref{cor:classification}; the group structure is the special case
of Theorem~\ref{thm:group-iff-normal}.
\end{remark}

\section{Reduction principles for decision}

\begin{definition}[Class function]
For a finite group $H$, a cost $J:H\to\R$ is a \emph{class function} if
$J(w x w^{-1}) = J(x)$ for all $w,x\in H$ (invariance under conjugation).
\end{definition}

\begin{proposition}[Class-function reduction]\label{prop:classfn}
If $L\trianglelefteq G$ (so $\Op\cong G/L$ is a group) and $J$ is a class
function on $\Op$, then $J$ is constant on each conjugacy class. Consequently
$\arg\min J$ is a union of conjugacy classes, and it suffices to evaluate $J$ on
one representative per class. The number of evaluations is the number of classes,
and the \emph{class equation}
\[
|\Op| = |Z(\Op)| + \sum_i [\,\Op : C_{\Op}(x_i)\,]
\]
enumerates the decision cells together with their multiplicities.
\end{proposition}

In the general (non-normal) case the relevant symmetry is the normalizer.

\begin{theorem}[Normalizer orbit decomposition]\label{thm:orbits}
The conjugation $\odot_{eg}\mapsto \odot_{e,wgw^{-1}}$ is well-defined on $\Op$
exactly for $w\in \Norm_G(L)$. The induced action of $\Norm_G(L)$ partitions
$\Op$ into \emph{decision cells} (orbits), counted by Burnside:
\[
\#\text{cells} \;=\; \frac{1}{|\Norm_G(L)|}\sum_{w\in \Norm_G(L)}
\bigl|\Fix_{\Op}(w)\bigr|.
\]
When $L\trianglelefteq G$ one has $\Norm_G(L)=G$ and the cells are exactly the
conjugacy classes of $\Op=G/L$, recovering Proposition~\ref{prop:classfn}.
\end{theorem}

\begin{proof}
For $w\in\Norm_G(L)$ and $a\in L$, $w(ag)w^{-1} = (waw^{-1})(wgw^{-1})$ with
$waw^{-1}\in L$, so $L\,w g w^{-1}$ depends only on $Lg$; the map is well-defined
on $\Op$ by Corollary~\ref{cor:classification}. For $w\notin\Norm_G(L)$ there is
$a\in L$ with $waw^{-1}\notin L$, and the image depends on the coset
representative. Orbit counting is Burnside's lemma; the normal case gives
$\Norm_G(L)=G$ acting on $G/L$ by conjugation, whose orbits are the conjugacy
classes.
\end{proof}

\begin{remark}
The quotient $\Norm_G(L)/L$ plays the role of a \emph{Weyl group} of the
configuration: it is the largest symmetry acting on $\Op$ through the kernel.
\end{remark}

\section{The class equation as a decision stratification}

Read as a spectrum of symmetry, the class equation stratifies the decision space:
\begin{itemize}
\item the \textbf{centre} $Z(\Op)$ consists of operations fixed by every inner
relabeling --- the frame-canonical, robust choices;
\item a \textbf{large class} ($[\Op:C(x)]$ large, small centralizer) is a
generic, low-symmetry operation;
\item a \textbf{small class} (large centralizer) is a highly symmetric,
structured operation.
\end{itemize}

\begin{remark}[Scope of structural reduction]
By the equivariance principle every twin operation is isomorphic to $\odot_e$;
hence a cost invariant under \emph{all} relabelings is constant on $\Op$ and
discriminates nothing. The class equation discriminates exactly the costs that
are invariant under the \emph{inner} symmetries (conjugation), such as the order
complexity $\kappa_{\mathrm{ord}}(\odot_{eg}) = \mathrm{ord}(gL)$. Frame-dependent
costs (conditioning in fixed coordinates, operation count) are handled by their
own symmetry group via Theorem~\ref{thm:orbits}.
\end{remark}

\section{Compact Lie groups: the Weyl integration formula}

For a compact connected Lie group $G$ with maximal torus $T$ and Weyl group $W$,
the finite class equation becomes the \emph{Weyl integration formula}: for a
class function $J$,
\[
\int_G J \, d\mu_{\mathrm{Haar}}
= \frac{1}{|W|}\int_T J(t)\,\bigl|\Delta(t)\bigr|^2\, dt,
\]
the continuous enumeration of conjugacy classes.

\begin{proposition}[$SU(2)$ instance]\label{prop:su2}
Conjugacy classes of $SU(2)$ are parametrized by the angle
$\theta\in[0,\pi]$ (eigenvalues $e^{\pm i\theta}$), and the push-forward of the
Haar measure is
\[
d\mu(\theta) = \frac{2}{\pi}\sin^2\theta \, d\theta, \qquad
\int_0^\pi d\mu = 1.
\]
A class cost $J(\theta)$ is decided by its Weyl average
$\int_0^\pi J(\theta)\,d\mu(\theta)$ and its optimum over $\theta\in[0,\pi]$.
\end{proposition}

This is the exact analogue of $\frac{1}{|G|}\sum_i |\text{class}_i|\,J(x_i)$:
the class equation, made continuous.

\begin{proposition}[Qudits: the $U(n)$/$SU(n)$ formula]\label{prop:sun}
For $U(n)$ the maximal torus is the diagonal subgroup, the Weyl group is $S_n$,
and the eigen-phase joint density is the squared Vandermonde
\[
d\mu(\theta_1,\dots,\theta_n)=\frac{1}{n!\,(2\pi)^n}
\prod_{j<k}\bigl|e^{i\theta_j}-e^{i\theta_k}\bigr|^2\,d\theta,
\]
the CUE eigenvalue density. The irreducible characters are the Schur
polynomials $\chi_\lambda=s_\lambda(e^{i\theta})$, and they are orthonormal,
$\langle\chi_\lambda,\chi_\mu\rangle=\delta_{\lambda\mu}$ --- the relation that
pins down the measure. A class cost on a qudit gate ($d$-level, under
conjugation by $SU(d)$) is a symmetric function of the eigen-phases, decided
against $d\mu$.
\end{proposition}

\begin{remark}[Universality persists]
A generic bilinear gate on $\mathbb{C}^d$ has a trivial structural kernel under
$SU(d)$, so $\Op\cong SU(d)$ for every $d$ tested --- the $SU(2)$ universality
instance extends to qudits.
\end{remark}

\section{Primary decomposition of the decision space (abelian case)}

When $G$ is finite abelian the kernel $L$ is automatically normal, so
$\Op=G/L$ is an abelian group and decomposes along its primary (Sylow-$p$)
components. Write $|G|=\prod_i p_i^{a_i}$ and let $G_p$ be the (unique) Sylow
$p$-subgroup.

\begin{theorem}[Primary decomposition]\label{thm:primary}
Let $G$ be finite abelian acting on $(E,\odot_e)$ with kernel $L$. Then
$G\cong\prod_p G_p$, every subgroup decomposes as $L=\prod_p L_p$ with
$L_p=L\cap G_p$ (so the compatibility hypothesis is automatic), and
\[
\Op = G/L \;\cong\; \prod_p\, G_p/L_p \;=\; \prod_p \Op_p .
\]
\end{theorem}

\begin{proof}
$G\cong\prod_p G_p$ is the primary decomposition of a finite abelian group; the
projection onto $G_p$ is $g\mapsto g^{e_p}$ where $e_p\equiv 1\ (\mathrm{mod}\ p^{a_p})$
and $e_p\equiv 0\ (\mathrm{mod}\ |G|/p^{a_p})$, whence $\prod_p g^{e_p}=g$. A
subgroup of $\prod_p G_p$ that is a product of $p$-groups is the product of its
intersections with the factors, giving $L=\prod_p L_p$. Quotienting factorwise
yields $G/L\cong\prod_p G_p/L_p$.
\end{proof}

\begin{corollary}[Unique factorisation of operations]\label{cor:crt}
Every twin operation factors uniquely as $\odot_{eg}\leftrightarrow (g_pL_p)_p
\in\prod_p\Op_p$, and $|\Op|=\prod_p|\Op_p|$.
\end{corollary}

\begin{proposition}[Separable decision]\label{prop:separable}
If the cost is separable, $J(\odot_{eg})=\sum_p J_p(g_pL_p)$, then
\[
\min_{\Op} J = \sum_p \min_{\Op_p} J_p,
\]
so the optimum is found componentwise. The search collapses from
$\prod_p|\Op_p|$ to $\sum_p|\Op_p|$ candidates --- the Chinese Remainder /
Fourier factorisation of the decision problem.
\end{proposition}

\begin{theorem}[Order-complexity partition]\label{thm:complexity}
Let $n=|\Op|$ and $C_d=\{\odot\in\Op:\kappa_{\mathrm{ord}}(\odot)=d\}$. Then
\[
\Op=\bigsqcup_{d\mid n} C_d, \qquad |C_d|=\varphi(d)\,m_d,
\]
where $m_d$ is the number of cyclic subgroups of $\Op$ of order $d$, and
$\sum_{d\mid n}|C_d|=n$.
\end{theorem}

\begin{proof}
Every $\odot\in\Op$ has an order $\kappa_{\mathrm{ord}}(\odot)=\mathrm{ord}(gL)$ dividing $n=|\Op|$ (Lagrange), so
$\Op=\bigsqcup_{d\mid n}C_d$ and $\sum_{d\mid n}|C_d|=n$. The elements of order $d$ are exactly the generators of
the cyclic subgroups of order $d$: each such subgroup has $\varphi(d)$ generators, and two order-$d$ elements
generate the same subgroup iff each is a power of the other. Summing over the $m_d$ cyclic subgroups of order $d$
gives $|C_d|=\varphi(d)\,m_d$.
\end{proof}

\section{Arithmetic-task costs (making the cost operational)}

The costs above were abstract. We now record \emph{measurable} costs tied to
genuine computation, taking $G=\langle\varphi\rangle$ a monogenic group so the
group element \emph{is} the computational domain. For a bijection $\varphi$ the
twin reduction of data $(x_i)$ is
$\bigoplus_\varphi x_i = \varphi^{-1}\!\bigl(\sum_i\varphi(x_i)\bigr)$, with
transformed total $T=\sum_i\varphi(x_i)$. For $\varphi=\log$ this is the
product, and $T$ is the log-likelihood.

Two well-posed decisions arise, each with a numeric cost $J$ that the tool
\emph{measures} rather than postulates:

\begin{itemize}
\item \textbf{Domain of computation} (fixed target $T$): evaluate $T$ in the
transformed domain ($\sum\varphi(x_i)$, stable) or the original domain
($\varphi(\bigoplus_\varphi x_i)$, prone to overflow/underflow). Cost: relative
error against a high-precision reference. For $\varphi=\log$ on wide-range data
the original domain underflows ($\prod x_i\to 0$, $\log 0=-\infty$) while the
transformed domain is exact to rounding.
\item \textbf{Optimal base} (representation): choose $\varphi$ so the data is
twin-linear; cost is the back-transformed fit RMSE. A power law selects
$\log$--$\log$, an exponential selects semi-$\log$, a line selects the identity.
\end{itemize}

\begin{remark}[Honest status]
These are \emph{not theorems} but the classical log-domain and linearisation
facts of numerical computing. The framework's contribution is to organise them
as a \emph{base-selection decision over $\langle\varphi\rangle$}, made decidable
by measurement; there is an honest speed/stability trade-off (the transformed
domain controls dynamic range at the price of more transcendental calls).
\end{remark}

\section{Computational realization}

The module \code{arithmetics/symbolic/twin/decision.py} implements these
statements; the test suite \code{arithmetics/tests/test\_twin\_decision.py}
verifies them. The correspondence is:

\begin{center}
\begin{tabular}{@{}lll@{}}
\toprule
\textbf{Statement} & \textbf{Code} & \textbf{Test} \\
\midrule
Class equation $|G|=|Z|+\sum[G:C]$ & \code{class\_equation} & \code{TestClassEquation} \\
Class-function reduction (Prop.~\ref{prop:classfn}) & \code{is\_class\_function}, \code{decision\_report} & \code{TestClassFunctions} \\
Classification $|\Op|=[G:L]$ (Cor.~\ref{cor:classification}) & \code{twin\_decision} & \code{test\_index\_equals\_num\_operations} \\
Group $\iff$ normal (Thm.~\ref{thm:group-iff-normal}, Rem.~\ref{rem:correction}) & \code{is\_normal}, \code{quotient} & \code{TestNormalityAndQuotient} \\
$S_4/{+}\bmod 4$ counterexample & \code{twin\_decision} & \code{test\_s4\_not\_normal\_weyl\_reduction} \\
Normalizer orbits (Thm.~\ref{thm:orbits}) & \code{normalizer}, \code{operation\_orbits} & \code{test\_orbits\_never\_exceed\_operations} \\
Weyl integration on $SU(2)$ (Prop.~\ref{prop:su2}) & \code{su2\_class\_average}, \code{su2\_class\_decision} & \code{TestWeylSU2} \\
$U(n)$/$SU(n)$ Weyl + characters (Prop.~\ref{prop:sun}) & \code{weyl\_average}, \code{verify\_character\_orthonormality} & \code{TestWeylMeasure}, \code{TestSchurCharacters} \\
Qudit universality (Rem.~after Prop.~\ref{prop:sun}) & \code{qudit\_universality} & \code{TestQuditDecisionAndUniversality} \\
Primary decomposition (Thm.~\ref{thm:primary}) & \code{verify\_primary\_decomposition} & \code{TestOperationDecomposition} \\
Unique factorisation (Cor.~\ref{cor:crt}) & \code{operation\_factorization} & \code{test\_unique\_factorization\_nonempty} \\
Separable decision (Prop.~\ref{prop:separable}) & \code{primary\_decision} & \code{TestSeparableDecision} \\
Complexity partition (Thm.~\ref{thm:complexity}) & \code{verify\_complexity\_partition} & \code{TestComplexityPartition} \\
Domain of computation (\S\,arith.) & \code{total\_error}, \code{recommend\_domain} & \code{TestTransformedTotal} \\
Optimal base (\S\,arith.) & \code{optimal\_base} & \code{TestOptimalBase} \\
\bottomrule
\end{tabular}
\end{center}

\section{Outlook}

The framework generalizes verbatim to abelian groups and direct products (where
$L$ is automatically normal and the class equation is the full enumeration),
to richer non-abelian groups (dihedral, symmetric, quaternion --- where the
class equation discriminates strongly), and to compact Lie groups via the Weyl
formula. The primary decomposition for abelian $G$
(Section~\ref{thm:primary}ff.) is realised; the remaining steps are:
(i)~cost functionals tied to genuine arithmetic tasks (conditioning, operation
count); and (ii)~the $SU(n)$ Weyl formula for the qudit instance.

\end{document}
